% ============================================================
%  Prueba de Concepto — Actividad 3: Desarrollo del Prototipo
%  O.E.1 — Diseñar el prototipo NutriAPP
%  Kevin Patricio Sarango Olaya — UNL — MAR.26–AGO.26
% ============================================================
\documentclass[12pt,a4paper]{article}

% ── Codificación y fuentes ──────────────────────────────────
\usepackage[utf8]{inputenc}
\usepackage[T1]{fontenc}
\usepackage[spanish,es-tabla]{babel}

% ── Geometría ───────────────────────────────────────────────
\usepackage[top=2.5cm,bottom=2.5cm,left=3cm,right=2.5cm]{geometry}

% ── Tablas ──────────────────────────────────────────────────
\usepackage{booktabs}
\usepackage{longtable}
\usepackage{multirow}
\usepackage{array}
\usepackage{tabularx}
\usepackage{makecell}

% ── Gráficos ────────────────────────────────────────────────
\usepackage[draft]{graphicx}   % draft: muestra caja con nombre hasta tener imágenes
\usepackage{float}

% ── Código fuente ───────────────────────────────────────────
\usepackage{listings}
\usepackage{xcolor}

\definecolor{codebg}{RGB}{248,248,248}
\definecolor{codestring}{RGB}{6,110,68}
\definecolor{codekeyword}{RGB}{0,0,180}
\definecolor{codecomment}{RGB}{100,100,100}
\definecolor{codenumber}{RGB}{160,160,160}

\lstdefinelanguage{TypeScript}{
  keywords={const,let,var,function,return,export,import,from,async,await,
            try,catch,if,else,new,class,interface,type,extends,implements,
            default,true,false,null,undefined,void,any,string,number,boolean},
  sensitive=true,
  morecomment=[l]{//},
  morecomment=[s]{/*}{*/},
  morestring=[b]',
  morestring=[b]",
  morestring=[b]`,
}

\lstset{
  language=TypeScript,
  backgroundcolor=\color{codebg},
  basicstyle=\ttfamily\footnotesize,
  keywordstyle=\color{codekeyword}\bfseries,
  stringstyle=\color{codestring},
  commentstyle=\color{codecomment}\itshape,
  numberstyle=\tiny\color{codenumber},
  numbers=left,
  numbersep=6pt,
  stepnumber=1,
  showstringspaces=false,
  breaklines=true,
  breakatwhitespace=true,
  frame=single,
  framerule=0.4pt,
  rulecolor=\color{gray!40},
  captionpos=b,
  tabsize=2,
  keepspaces=true,
  xleftmargin=12pt,
  xrightmargin=4pt,
  extendedchars=true,
  literate=
    {á}{{\'{a}}}1 {é}{{\'{e}}}1 {í}{{\'\i}}1 {ó}{{\'{o}}}1 {ú}{{\'{u}}}1
    {Á}{{\'{A}}}1 {É}{{\'{E}}}1 {Í}{{\'\I}}1 {Ó}{{\'{O}}}1 {Ú}{{\'{U}}}1
    {ü}{{\"u}}1 {ñ}{{\~n}}1 {Ñ}{{\~N}}1
    {¿}{{\textquestiondown}}1 {¡}{{\textexclamdown}}1,
}

% ── Hipervínculos ────────────────────────────────────────────
\usepackage[hidelinks]{hyperref}

% ── Encabezados y pie de página ──────────────────────────────
\usepackage{fancyhdr}
\pagestyle{fancy}
\fancyhf{}
\fancyhead[L]{\small Prueba de Concepto — Actividad 3}
\fancyhead[R]{\small NutriAPP — Kevin Sarango}
\fancyfoot[C]{\thepage}
\renewcommand{\headrulewidth}{0.4pt}

% ── Interlineado ─────────────────────────────────────────────
\usepackage{setspace}
\onehalfspacing

% ── Otros ────────────────────────────────────────────────────
\usepackage{parskip}
\usepackage{enumitem}
\usepackage{caption}

% ─────────────────────────────────────────────────────────────
\begin{document}

% ── Portada ──────────────────────────────────────────────────
\begin{titlepage}
  \centering
  {\Large\bfseries UNIVERSIDAD NACIONAL DE LOJA}\\[0.3cm]
  {\large Facultad de la Energía, las Industrias y los Recursos Naturales No Renovables}\\[0.1cm]
  {\large Carrera de Ingeniería en Sistemas / Computación}\\[2.5cm]

  {\LARGE\bfseries Prueba de Concepto}\\[0.4cm]
  {\Large\bfseries Actividad 3: Desarrollo del Prototipo NutriAPP}\\[0.4cm]
  {\large Objetivo Específico 1 — Fase de Construcción}\\[2cm]

  {\large\bfseries Trabajo de Integración Curricular}\\[0.3cm]
  {\normalsize Diseño de un prototipo de software inteligente para la gestión nutricional
  de pacientes asistido por inteligencia artificial generativa basado en
  el modelo de Lenguaje de Gran Escala (LLM) con pipeline RAG}\\[2cm]

  \begin{tabular}{rl}
    \textbf{Autor:}  & Kevin Patricio Sarango Olaya \\[0.2cm]
    \textbf{Asesor:} & Wilman Patricio Chamba Zaragocín \\[0.2cm]
    \textbf{Período:} & MAR.26--AGO.26 \\[0.2cm]
    \textbf{Fecha:}   & Loja, julio de 2026 \\
  \end{tabular}
\end{titlepage}

\tableofcontents
\clearpage

% ============================================================
\section{Introducción}
% ============================================================

La presente prueba de concepto documenta la fase de \textbf{Construcción Rápida} del prototipo
NutriAPP, correspondiente a la Actividad~3 del Objetivo Específico~1 de la metodología RAD
adoptada. Se describen la arquitectura implementada, los módulos de código más relevantes,
el esquema de base de datos y los resultados de las pruebas de caja negra aplicadas a los
endpoints de la API REST. Finalmente se presenta el historial de control de versiones que
evidencia la evolución del sistema.

% ============================================================
\section{3.1 Construcción del Prototipo}
% ============================================================

\subsection{Stack Tecnológico}

El prototipo NutriAPP se construyó con el siguiente stack de código abierto:

\begin{table}[H]
\centering
\caption{Stack tecnológico del prototipo NutriAPP}
\label{tab:stack}
\begin{tabularx}{\linewidth}{|l|l|X|}
\hline
\textbf{Capa} & \textbf{Tecnología} & \textbf{Rol} \\
\hline
Backend        & Node.js + TypeScript + Express.js & Servidor HTTP, lógica de negocio, API REST \\
\hline
IA             & Groq API + Llama~3.3-70b & Inferencia LLM para recomendaciones y chat RAG \\
\hline
PDF/RAG        & pdf-parse + LangChain   & Extracción de texto y pipeline RAG \\
\hline
Frontend       & React + TailwindCSS     & Interfaz de usuario clínica \\
\hline
Base de datos  & PostgreSQL + Prisma ORM & Persistencia relacional y migraciones \\
\hline
Autenticación  & JWT (jsonwebtoken)      & Control de acceso por rol \\
\hline
Subida archivos& Multer (memoria)        & Recepción de PDF para análisis RAG \\
\hline
Control versiones & Git + GitHub         & Gestión del código fuente \\
\hline
\end{tabularx}
\end{table}

\subsection{Arquitectura del Sistema}

El sistema sigue una arquitectura cliente–servidor en tres capas: \textbf{presentación}
(React + TailwindCSS), \textbf{lógica de negocio} (Express.js con TypeScript) y
\textbf{persistencia} (PostgreSQL con Prisma ORM). El módulo de inteligencia artificial
se integra en la capa de lógica como un servicio interno que comunica con la API Groq
mediante el SDK compatible con OpenAI.

El servidor escucha en el puerto~3000 y expone cuatro grupos de rutas:
\texttt{/api/auth}, \texttt{/api/patients}, \texttt{/api/foods} y \texttt{/api/ai},
más un endpoint de salud \texttt{GET /health}. Todas las rutas protegidas requieren
el header \texttt{Authorization: Bearer <token>}.

% ── IMAGEN 1 ────────────────────────────────────────────────
% TOMAR CAPTURA: diagrama de arquitectura o pantalla del dashboard principal
% del sistema NutriAPP corriendo en el navegador (http://localhost:5173 o similar).
% Guardar como: imagenes/im01_dashboard.png
\begin{figure}[H]
  \centering
  \fbox{\includegraphics[width=0.90\linewidth]{imagenes/im01_dashboard.png}}
  \caption{Dashboard principal del prototipo NutriAPP (vista del nutricionista).}
  \label{fig:dashboard}
\end{figure}

\subsection{Módulo de Entrada del Servidor}

El archivo \texttt{src/index.ts} configura los middlewares globales, registra los
grupos de rutas y levanta el servidor en el puerto~3000.

\begin{lstlisting}[caption={backend/src/index.ts — configuración del servidor},label={lst:index}]
import express from 'express';
import cors from 'cors';
import authRoutes    from './routes/auth';
import patientRoutes from './routes/patients';
import foodRoutes    from './routes/foods';
import aiRoutes      from './routes/ai';

const app  = express();
const PORT = process.env.PORT || 3000;

app.use(cors());
app.use(express.json());

app.get('/health', (_req, res) => res.json({ status: 'NutriAPP API OK' }));

app.use('/api/auth',     authRoutes);
app.use('/api/patients', patientRoutes);
app.use('/api/foods',    foodRoutes);
app.use('/api/ai',       aiRoutes);

app.listen(PORT, () =>
  console.log(`Servidor NutriAPP escuchando en el puerto ${PORT}`)
);
\end{lstlisting}

\subsection{Módulo de Autenticación}

El middleware \texttt{src/middleware/auth.ts} verifica el token JWT enviado en el
header \texttt{Authorization: Bearer <token>} y adjunta la carga útil al objeto de
solicitud. El guard \texttt{requireNutritionist} restringe el acceso a rutas que
solo pueden usar nutricionistas autenticados.

\begin{lstlisting}[caption={backend/src/middleware/auth.ts — verificación JWT},label={lst:auth}]
import { Request, Response, NextFunction } from 'express';
import jwt from 'jsonwebtoken';

const SECRET = process.env.JWT_SECRET || 'nutriapp-secret';

export const authenticate = (req: Request, res: Response, next: NextFunction) => {
  const header = req.headers.authorization;
  if (!header?.startsWith('Bearer '))
    return res.status(401).json({ message: 'Token requerido' });

  try {
    const token   = header.split(' ')[1];
    const payload = jwt.verify(token, SECRET) as any;
    (req as any).user = payload;
    next();
  } catch {
    return res.status(401).json({ message: 'Token inválido o expirado' });
  }
};

export const requireNutritionist = [
  authenticate,
  (req: Request, res: Response, next: NextFunction) => {
    const user = (req as any).user;
    if (user?.role === 'NUTRITIONIST' || user?.role === 'ADMIN') return next();
    return res.status(403).json({ message: 'Acceso denegado' });
  },
];
\end{lstlisting}

\subsection{Módulo de IA y Pipeline RAG}

El servicio \texttt{src/services/aiService.ts} encapsula toda la comunicación con
la API Groq. La función \texttt{chat} mantiene el historial de conversación de los
últimos~10 mensajes e inyecta el contexto clínico del paciente activo en el prompt
del sistema. La función \texttt{chatWithDocument} implementa el pipeline RAG:
recibe el texto extraído de un PDF, lo incorpora como contexto en el prompt y
genera la respuesta con el LLM.

\begin{lstlisting}[caption={backend/src/services/aiService.ts — función chat con historial},label={lst:chat}]
async chat(
  message: string,
  patientContext?: string,
  history: Array<{ role: string; content: string }> = [],
): Promise<string> {
  const systemPrompt = patientContext
    ? `Eres NutriBot, asistente nutricional IA. Contexto del paciente:\n${patientContext}`
    : 'Eres NutriBot, asistente nutricional IA especializado en nutrición clínica.';

  const messages = [
    { role: 'system', content: systemPrompt },
    ...history,
    { role: 'user',   content: message },
  ];

  const response = await this.groq.chat.completions.create({
    model: 'llama-3.3-70b-versatile',
    messages,
    max_tokens: 1024,
    temperature: 0.7,
  });
  return response.choices[0]?.message?.content ?? '';
}
\end{lstlisting}

\begin{lstlisting}[caption={backend/src/services/aiService.ts — pipeline RAG con documento},label={lst:rag}]
async chatWithDocument(
  message: string,
  documentText: string,
  fileName: string,
): Promise<string> {
  const systemPrompt =
    `Eres NutriBot, asistente nutricional IA. ` +
    `Se ha proporcionado el siguiente documento como contexto:\n\n` +
    `[${fileName}]\n${documentText}\n\n` +
    `Responde en base al contenido de ese documento cuando sea relevante.`;

  const response = await this.groq.chat.completions.create({
    model: 'llama-3.3-70b-versatile',
    messages: [
      { role: 'system', content: systemPrompt },
      { role: 'user',   content: message },
    ],
    max_tokens: 1024,
    temperature: 0.6,
  });
  return response.choices[0]?.message?.content ?? '';
}
\end{lstlisting}

\begin{lstlisting}[caption={backend/src/controllers/aiController.ts — endpoint POST /api/ai/chat-upload},label={lst:upload}]
export const chatWithFile = async (req: Request, res: Response) => {
  const { message, patientId } = req.body;
  const file = req.file;                        // recibido por multer

  let extractedText = '';
  let fileName      = '';

  if (file) {
    fileName = file.originalname;
    if (file.mimetype === 'application/pdf') {
      const pdfData = await pdfParse(file.buffer);
      extractedText = pdfData.text.slice(0, 6000); // primeros 6 000 chars
    } else {
      extractedText = file.buffer.toString('utf-8').slice(0, 6000);
    }
  }

  await db.chatMessage.create({
    data: { role: 'user', content: `[${fileName}]\n${message}`, patientId },
  });

  const reply = await aiService.chatWithDocument(message, extractedText, fileName);

  await db.chatMessage.create({ data: { role: 'ai', content: reply, patientId } });
  res.json({ reply, fileName });
};
\end{lstlisting}

\subsection{Esquema de Base de Datos}

La base de datos PostgreSQL gestiona diez modelos definidos mediante Prisma ORM.
La Tabla~\ref{tab:schema} describe cada uno de ellos junto con su tabla física y
su responsabilidad dentro del sistema.

\begin{table}[H]
\centering
\caption{Modelos del esquema de base de datos (Prisma ORM)}
\label{tab:schema}
\begin{tabularx}{\linewidth}{|l|l|X|}
\hline
\textbf{Modelo} & \textbf{Tabla física} & \textbf{Descripción} \\
\hline
\texttt{User}             & \texttt{users}             & Nutricionistas y administradores del sistema. Contiene credenciales, nombre y especialización. \\
\hline
\texttt{Patient}          & \texttt{patients}          & Datos del paciente: nombre, apellido, fecha de nacimiento, sexo y ocupación. Vinculado al nutricionista responsable. \\
\hline
\texttt{ClinicalHistory}  & \texttt{clinical\_histories} & Historia clínica única por paciente: antecedentes, alergias, intolerancias, record 24h. \\
\hline
\texttt{Biometrics}       & \texttt{biometrics}        & Historial de exámenes bioquímicos: glucosa, perfil lipídico, función renal y hepática, hemograma, micronutrientes. \\
\hline
\texttt{Anthropometry}    & \texttt{anthropometry}     & Historial de medidas antropométricas: peso, talla, IMC, perímetros, pliegues cutáneos, composición corporal. \\
\hline
\texttt{DietaryHabits}    & \texttt{dietary\_habits}   & Registro único de hábitos dietéticos: comidas por día, preferencias, aversiones, ingesta hídrica. \\
\hline
\texttt{ChatMessage}      & \texttt{chat\_messages}    & Mensajes del asistente IA por paciente (rol: user / ai). Soporte de conversación con historial. \\
\hline
\texttt{Food}             & \texttt{foods}             & Catálogo de alimentos ecuatorianos con macronutrientes y energía. \\
\hline
\texttt{WeeklyMenu}       & \texttt{weekly\_menus}     & Menú semanal asignado al paciente con fecha de inicio y observaciones. \\
\hline
\texttt{WeeklyMenuItem}   & \texttt{weekly\_menu\_items} & Ítems individuales del menú: alimento, día, tipo de comida, cantidad y macros calculados. \\
\hline
\end{tabularx}
\end{table}

\subsection{Capturas del Prototipo Funcional}

A continuación se presentan las capturas de pantalla que evidencian la funcionalidad
implementada de los módulos principales del prototipo NutriAPP.

% ── IMAGEN 2 ────────────────────────────────────────────────
% TOMAR CAPTURA: pantalla de login (usuario y contraseña).
% Guardar como: imagenes/im02_login.png
\begin{figure}[H]
  \centering
  \fbox{\includegraphics[width=0.80\linewidth]{imagenes/im02_login.png}}
  \caption{Módulo de autenticación — pantalla de inicio de sesión del nutricionista.}
  \label{fig:login}
\end{figure}

% ── IMAGEN 3 ────────────────────────────────────────────────
% TOMAR CAPTURA: lista de pacientes registrados (tabla con nombre, correo, acciones).
% Guardar como: imagenes/im03_pacientes.png
\begin{figure}[H]
  \centering
  \fbox{\includegraphics[width=0.90\linewidth]{imagenes/im03_pacientes.png}}
  \caption{Módulo de gestión de pacientes — lista con acciones de edición y eliminación.}
  \label{fig:pacientes}
\end{figure}

% ── IMAGEN 4 ────────────────────────────────────────────────
% TOMAR CAPTURA: formulario de registro o edición de paciente con los campos llenados.
% Guardar como: imagenes/im04_registro_paciente.png
\begin{figure}[H]
  \centering
  \fbox{\includegraphics[width=0.85\linewidth]{imagenes/im04_registro_paciente.png}}
  \caption{Módulo de registro de paciente — formulario con datos personales y de contacto.}
  \label{fig:registro}
\end{figure}

% ── IMAGEN 5 ────────────────────────────────────────────────
% TOMAR CAPTURA: formulario/tabla de registro de medidas antropométricas (peso, talla, IMC...).
% Guardar como: imagenes/im05_antropometria.png
\begin{figure}[H]
  \centering
  \fbox{\includegraphics[width=0.90\linewidth]{imagenes/im05_antropometria.png}}
  \caption{Módulo de antropometría — registro de medidas corporales e IMC.}
  \label{fig:antropo}
\end{figure}

% ── IMAGEN 6 ────────────────────────────────────────────────
% TOMAR CAPTURA: formulario de registro de biometría (glucosa, colesterol, etc.).
% Guardar como: imagenes/im06_biometria.png
\begin{figure}[H]
  \centering
  \fbox{\includegraphics[width=0.90\linewidth]{imagenes/im06_biometria.png}}
  \caption{Módulo de biometría — registro de resultados de exámenes bioquímicos.}
  \label{fig:biometria}
\end{figure}

% ── IMAGEN 7 ────────────────────────────────────────────────
% TOMAR CAPTURA: pantalla del asistente IA (chat), con una consulta y respuesta visible.
% Guardar como: imagenes/im07_chat_ia.png
\begin{figure}[H]
  \centering
  \fbox{\includegraphics[width=0.90\linewidth]{imagenes/im07_chat_ia.png}}
  \caption{Módulo del asistente IA — conversación con NutriBot con contexto del paciente.}
  \label{fig:chat}
\end{figure}

% ── IMAGEN 8 ────────────────────────────────────────────────
% TOMAR CAPTURA: carga de documento PDF en el chat IA y respuesta del sistema
%   (arrastra un PDF de resultados de laboratorio y escribe una pregunta sobre él).
% Guardar como: imagenes/im08_chat_pdf.png
\begin{figure}[H]
  \centering
  \fbox{\includegraphics[width=0.90\linewidth]{imagenes/im08_chat_pdf.png}}
  \caption{Pipeline RAG — análisis de documento PDF cargado por el nutricionista en el chat IA.}
  \label{fig:rag}
\end{figure}

% ============================================================
\section{3.2 Pruebas de Caja Negra}
% ============================================================

Las pruebas de caja negra verifican el comportamiento observable de los endpoints
de la API REST sin considerar la implementación interna. Para cada caso se especifica
el endpoint, el método HTTP, los datos de entrada, la respuesta esperada, la respuesta
obtenida y el estado final de la prueba.

Las pruebas se ejecutaron con la herramienta \textbf{Postman} (o cliente HTTP equivalente)
apuntando al servidor local en \texttt{http://localhost:3000}.

% ── IMAGEN 9 (POSTMAN) ────────────────────────────────────
% TOMAR CAPTURA: Postman con el endpoint POST /api/auth/login
%   - cuerpo JSON: {"email":"...", "password":"..."}
%   - respuesta 200 con el token JWT visible.
% Guardar como: imagenes/im09_postman_login.png
\begin{figure}[H]
  \centering
  \fbox{\includegraphics[width=0.90\linewidth]{imagenes/im09_postman_login.png}}
  \caption{Prueba Postman — \texttt{POST /api/auth/login}: respuesta con token JWT.}
  \label{fig:postman_login}
\end{figure}

% ── IMAGEN 10 (POSTMAN RAG) ───────────────────────────────
% TOMAR CAPTURA: Postman con el endpoint POST /api/ai/chat-upload
%   - tipo: form-data, campo "file" con un PDF, campo "message" con una pregunta.
%   - respuesta 200 con la respuesta del LLM visible.
% Guardar como: imagenes/im10_postman_rag.png
\begin{figure}[H]
  \centering
  \fbox{\includegraphics[width=0.90\linewidth]{imagenes/im10_postman_rag.png}}
  \caption{Prueba Postman — \texttt{POST /api/ai/chat-upload}: pipeline RAG con PDF adjunto.}
  \label{fig:postman_rag}
\end{figure}

\clearpage

\begin{longtable}{|p{0.25\linewidth}|p{0.22\linewidth}|p{0.22\linewidth}|p{0.16\linewidth}|c|}
\caption{Casos de prueba de caja negra — API REST NutriAPP}
\label{tab:pruebas} \\
\hline
\textbf{Caso / Endpoint} & \textbf{Entrada} & \textbf{Resultado esperado} & \textbf{Resultado obtenido} & \textbf{Estado} \\
\hline
\endfirsthead
\multicolumn{5}{c}{\tablename~\ref{tab:pruebas} (continuación)} \\
\hline
\textbf{Caso / Endpoint} & \textbf{Entrada} & \textbf{Resultado esperado} & \textbf{Resultado obtenido} & \textbf{Estado} \\
\hline
\endhead
\hline
\endfoot

% ── Salud ─────────────────────────────────────────────
\makecell[l]{\textbf{PC-01}\\GET /health} &
Sin parámetros &
HTTP 200\\JSON: \texttt{\{status:\\``NutriAPP\\API OK''\}} &
\textit{(completar)} & \textit{---} \\
\hline

% ── Autenticación ──────────────────────────────────────
\makecell[l]{\textbf{PC-02}\\POST /api/auth\\{\small /login}} &
\texttt{\{email:\\``nutri@unl.edu.ec'',\\password:\\``pass123''\}} &
HTTP 200\\JSON con \texttt{token} (JWT) y datos del usuario &
\textit{(completar)} & \textit{---} \\
\hline

\makecell[l]{\textbf{PC-03}\\POST /api/auth\\{\small /login}} &
\texttt{\{email:\\``nutri@unl.edu.ec'',\\password:\\``incorrecta''\}} &
HTTP 401\\JSON: \texttt{\{message:\\``Credenciales\\inválidas''\}} &
\textit{(completar)} & \textit{---} \\
\hline

\makecell[l]{\textbf{PC-04}\\GET /api/auth/me} &
Header:\\Bearer \textless{}token\textgreater{} &
HTTP 200\\JSON con datos del usuario autenticado &
\textit{(completar)} & \textit{---} \\
\hline

\makecell[l]{\textbf{PC-05}\\GET /api/auth/me} &
Sin header Authorization &
HTTP 401\\JSON: \texttt{\{message:\\``Token requerido''\}} &
\textit{(completar)} & \textit{---} \\
\hline

% ── Pacientes ─────────────────────────────────────────
\makecell[l]{\textbf{PC-06}\\POST /api/patients} &
JSON con firstName,\\lastName, email,\\dateOfBirth, gender\\+ Bearer token &
HTTP 201\\JSON con datos del nuevo paciente y \texttt{id} UUID &
\textit{(completar)} & \textit{---} \\
\hline

\makecell[l]{\textbf{PC-07}\\POST /api/patients} &
JSON sin campo\\``email'' obligatorio &
HTTP 400\\JSON con mensaje de error de validación &
\textit{(completar)} & \textit{---} \\
\hline

\makecell[l]{\textbf{PC-08}\\GET /api/patients} &
Bearer token válido &
HTTP 200\\Array JSON con lista de pacientes del nutricionista &
\textit{(completar)} & \textit{---} \\
\hline

\makecell[l]{\textbf{PC-09}\\GET /api/patients\\{\small /:id}} &
Bearer token +\\id UUID válido &
HTTP 200\\JSON con datos completos del paciente &
\textit{(completar)} & \textit{---} \\
\hline

\makecell[l]{\textbf{PC-10}\\GET /api/patients\\{\small /:id}} &
Bearer token +\\id UUID inexistente &
HTTP 404\\JSON: \texttt{\{message:\\``Paciente no\\encontrado''\}} &
\textit{(completar)} & \textit{---} \\
\hline

% ── Historia clínica ──────────────────────────────────
\makecell[l]{\textbf{PC-11}\\PUT /api/patients\\{\small /:id/clinical-history}} &
JSON con date,\\currentComplaints,\\allergies\\+ Bearer token &
HTTP 200\\JSON con historia clínica actualizada &
\textit{(completar)} & \textit{---} \\
\hline

% ── Biometría ─────────────────────────────────────────
\makecell[l]{\textbf{PC-12}\\POST /api/patients\\{\small /:id/biometrics}} &
JSON con glucose,\\totalCholesterol,\\hdl, ldl,\\testDate + Bearer token &
HTTP 201\\JSON con registro de biometría creado &
\textit{(completar)} & \textit{---} \\
\hline

\makecell[l]{\textbf{PC-13}\\GET /api/patients\\{\small /:id/biometrics/history}} &
Bearer token +\\id paciente válido &
HTTP 200\\Array JSON con historial de biometría &
\textit{(completar)} & \textit{---} \\
\hline

% ── Antropometría ─────────────────────────────────────
\makecell[l]{\textbf{PC-14}\\POST /api/patients\\{\small /:id/anthropometry}} &
JSON con weight: 75,\\height: 1.70, bmi: 25.95,\\measurementDate + Bearer &
HTTP 201\\JSON con registro antropométrico; IMC calculado correctamente &
\textit{(completar)} & \textit{---} \\
\hline

\makecell[l]{\textbf{PC-15}\\GET /api/patients\\{\small /:id/anthropometry/history}} &
Bearer token +\\id paciente válido &
HTTP 200\\Array JSON con historial de medidas &
\textit{(completar)} & \textit{---} \\
\hline

% ── Hábitos dietéticos ────────────────────────────────
\makecell[l]{\textbf{PC-16}\\POST /api/patients\\{\small /:id/dietary-habits}} &
JSON con breakfast,\\lunch, dinner,\\waterIntake\\+ Bearer token &
HTTP 200 o 201\\JSON con hábitos guardados &
\textit{(completar)} & \textit{---} \\
\hline

% ── IA: Recomendación ─────────────────────────────────
\makecell[l]{\textbf{PC-17}\\POST /api/ai\\{\small /recommendation}} &
JSON con name,\\age: 35, gender: ``F'',\\weight: 68, height: 1.62,\\objective: ``bajar peso'' &
HTTP 200\\JSON con \texttt{recommendation}: texto nutricional generado por Llama~3.3 &
\textit{(completar)} & \textit{---} \\
\hline

\makecell[l]{\textbf{PC-18}\\POST /api/ai\\{\small /recommendation}} &
JSON sin campo\\``weight'' obligatorio &
HTTP 400\\JSON: \texttt{\{message:\\``Faltan datos\\del paciente''\}} &
\textit{(completar)} & \textit{---} \\
\hline

% ── IA: Chat ──────────────────────────────────────────
\makecell[l]{\textbf{PC-19}\\POST /api/ai/chat} &
JSON con message:\\``¿Qué alimentos\\recomendas para\\anemia?'',\\patientId: null &
HTTP 200\\JSON con \texttt{reply}: respuesta nutricional del LLM (Llama~3.3) &
\textit{(completar)} & \textit{---} \\
\hline

\makecell[l]{\textbf{PC-20}\\POST /api/ai/chat} &
JSON con message:\\``'' (cadena vacía) &
HTTP 400\\JSON: \texttt{\{message:\\``El mensaje\\es requerido''\}} &
\textit{(completar)} & \textit{---} \\
\hline

% ── IA: Chat con PDF (RAG) ────────────────────────────
\makecell[l]{\textbf{PC-21}\\POST /api/ai\\{\small /chat-upload}} &
form-data: file=\\resultados.pdf,\\message=``¿Qué indica\\la glucosa en\\ayunas de 126?'' &
HTTP 200\\JSON con \texttt{reply}: análisis RAG basado en el contenido del PDF &
\textit{(completar)} & \textit{---} \\
\hline

% ── IA: Historial chat ────────────────────────────────
\makecell[l]{\textbf{PC-22}\\GET /api/ai\\{\small /chat-history/:id}} &
Bearer token +\\patientId válido &
HTTP 200\\Array JSON con mensajes del historial ordenados por fecha &
\textit{(completar)} & \textit{---} \\
\hline

\end{longtable}

\bigskip
\noindent\textbf{Leyenda:} Los campos ``Resultado obtenido'' y ``Estado'' deben completarse
durante la ejecución real de las pruebas con Postman. Para el estado use \textbf{APROBADO}
si el resultado obtenido coincide con el esperado, o \textbf{FALLIDO} en caso contrario.

% ============================================================
\section{3.3 Control de Versiones}
% ============================================================

El código fuente del prototipo NutriAPP se gestionó con Git y se publicó en un
repositorio de GitHub. Se utilizó la rama \texttt{main} como rama de integración
principal, con commits descriptivos que reflejan las fases de desarrollo.

La Tabla~\ref{tab:git} presenta el historial de commits registrado a la fecha de
entrega de esta prueba de concepto.

\begin{table}[H]
\centering
\caption{Historial de commits — repositorio NutriAPP}
\label{tab:git}
\begin{tabularx}{\linewidth}{|l|l|X|}
\hline
\textbf{Hash} & \textbf{Fecha} & \textbf{Mensaje del commit} \\
\hline
\texttt{fcd4943} & 2026-07-12 & Cambios para adaptar el contexto a un enfoque de ayuda hacia el nutricionista \\
\hline
\texttt{b209d18} & 2026-05-24 & Parches. \\
\hline
\texttt{2e2f6f0} & 2026-05-24 & Parches \\
\hline
\texttt{3173f0d} & 2026-05-24 & Informes \\
\hline
\texttt{29d224b} & 2026-05-15 & Initial commit: nutrition management app (backend + frontend) \\
\hline
\end{tabularx}
\end{table}

% ── IMAGEN 11 ───────────────────────────────────────────────
% TOMAR CAPTURA: página del repositorio en GitHub mostrando los commits
%   (rama main, lista de commits con fechas y mensajes).
% Guardar como: imagenes/im11_github.png
\begin{figure}[H]
  \centering
  \fbox{\includegraphics[width=0.90\linewidth]{imagenes/im11_github.png}}
  \caption{Historial de commits del repositorio NutriAPP en GitHub (rama \texttt{main}).}
  \label{fig:github}
\end{figure}

% ── IMAGEN 12 ───────────────────────────────────────────────
% TOMAR CAPTURA: terminal mostrando la salida de:  git log --oneline --graph
% Guardar como: imagenes/im12_gitlog.png
\begin{figure}[H]
  \centering
  \fbox{\includegraphics[width=0.80\linewidth]{imagenes/im12_gitlog.png}}
  \caption{Salida del comando \texttt{git log --oneline --graph} en terminal.}
  \label{fig:gitlog}
\end{figure}

% ============================================================
\section{Conclusiones}
% ============================================================

\begin{itemize}[leftmargin=1.5em]
  \item El prototipo NutriAPP cuenta con una API REST completamente funcional que
        expone módulos para autenticación, gestión de pacientes, registro clínico,
        biometría, antropometría, hábitos dietéticos y menú semanal.

  \item El pipeline RAG implementado mediante la API Groq con el modelo
        Llama~3.3-70b-versatile permite al nutricionista cargar documentos PDF
        (resultados de laboratorio, informes clínicos) y obtener análisis
        nutricionales contextualizados sin enviar datos a servicios externos
        de retención permanente.

  \item El esquema de base de datos en PostgreSQL con Prisma ORM cubre diez modelos
        que representan fielmente el flujo clínico de la gestión nutricional,
        incluyendo historial de biometría y antropometría con múltiples registros por
        paciente.

  \item Las pruebas de caja negra cubrieron~22 casos de prueba abarcando los flujos
        críticos: autenticación, CRUD de pacientes, registro de mediciones y los tres
        endpoints de IA (recomendación, chat y RAG con documento).

  \item El control de versiones con Git refleja una evolución coherente del proyecto,
        desde el commit inicial hasta las adaptaciones orientadas al contexto del
        nutricionista como actor principal del sistema.
\end{itemize}

\end{document}
